% ============================================================
%  comment_response.tex
%  Point-by-point response to the supervisor's comments on
%  "Multi-server queueing with reinforcement learning"
%  (version read up to Section 7.1).
%
%  Written in English so it can be reused directly for the paper
%  revision; the original comments are quoted verbatim in Dutch.
% ============================================================
\documentclass[11pt,a4paper]{article}

\usepackage{amsmath,amssymb,amsthm}
\usepackage{geometry}
\geometry{margin=2.5cm}
\usepackage{booktabs}
\usepackage{array}
\usepackage{hyperref}
\usepackage{xcolor}
\usepackage{microtype}
\usepackage{parskip}

\newcommand{\E}{\mathbb{E}}
\newcommand{\Prob}{\mathbb{P}}
\newcommand{\comment}[1]{\textcolor{red!70!black}{\textit{``#1''}}}

\title{Response to comments on\\\emph{Multi-server queueing with reinforcement learning}}
\author{}
\date{}

\begin{document}
\maketitle

\begin{abstract}
\noindent
Point-by-point response to the comments on the manuscript version read up to
Section~7.1.  The comments are quoted verbatim; responses give full
derivations where a mathematical claim is at stake.  Three groups are
distinguished: (A) genuine errors, with corrections; (B) modelling questions
that have since been answered by experiments --- in one case the comment
anticipates exactly the reformulation that was subsequently implemented; and
(C) clarifications where the manuscript is correct but underexplained.
Terminology and notation points are collected in (D).
\end{abstract}

% ============================================================
\section*{A. Errors, with corrections}
% ============================================================

\subsection*{A.1 The arrival term in the exit rate is inverted}

\comment{klopt de 1e term wel? bij aankomst krijg je alleen een verandering in
toestand als W $-1$ is, dan gaat of W van 0 of een server gaat aan het werk}
\hfill (\S2.3)

\textbf{The comment is correct; the indicator is inverted.}  The manuscript
states
\[
  \Lambda(x) \;=\; \sum_{n=1}^{N}\lambda_n \mathbf{1}\{W_n^1(x) > 0\}
  \;+\; \sum_{k=1}^{K}\sum_{m\in G_k} y_{m,k}(x)\,\mu(m,k) \;+\; \gamma .
\]

In the FIL representation only one waiting job per type is tracked.  A type-$n$
arrival therefore changes the state \emph{only} when queue $n$ is empty; if a
job is already tracked ($W_n^1 \ge 0$), a further arrival is not represented in
the state at all --- it is absorbed into the conditional model of the untracked
queue and revealed later by the FIL-refresh distribution~(3).  Consequently the
arrival stream of type $n$ contributes to the exit rate exactly when
$W_n^1 = -1$:
\[
  \boxed{\;\Lambda(x) \;=\; \sum_{n=1}^{N}\lambda_n \mathbf{1}\{W_n^1(x) = -1\}
  \;+\; \sum_{k=1}^{K}\sum_{m\in G_k} y_{m,k}(x)\,\mu(m,k) \;+\; \gamma\;}
\]
The accompanying sentence in \S3.2 (``if level $W_n^1 > 0$ \ldots arrivals do no
longer take place'') should likewise read $W_n^1 \ge 0$.  The implementation
uses the correct condition, so no numerical result is affected: this is an
error in the write-up only.

\subsection*{A.2 The M/M/1 validation cannot distinguish the two objectives}

\comment{is g niet lambda * rho? komt op hetzelfde neer omdat mu=1}
\hfill (\S6.4, Table 2)

\comment{je moet toch vanaf 1 sommeren? dan krijg je g = rho. en g = lambda *
P(vertraging)} \hfill (\S6.3)

These two comments are one observation, and it identifies a real weakness in
the validation experiment.  Both derivations below are correct; they describe
\emph{different objectives} that happen to coincide numerically at $\mu = 1$.

\paragraph{Lateness (what the manuscript computes).}
Cost accrues continuously while a job is overdue.  With $D = 0$, cost $c_1/\gamma$
is charged at each tick for which the FIL is overdue, and ticks occur at rate
$\gamma$, so the cost per unit real time is $c_1\,\Prob(W^1 > 0)$.

Two points deserve stating explicitly, because both were queried.  First, a job
\emph{in service} has left the queue and is no longer the FIL; ``a job is
waiting'' therefore requires $n \ge 2$ in M/M/1 (one in service, at least one
queued), which is why the sum starts at $n = 2$ and not at $n = 1$.  Second,
the tick \emph{first} advances all waiting-time levels and cost is charged
\emph{afterwards}; a job present at the tick with level $0$ becomes level $1$
and is charged.  Hence every queued job at a tick instant is charged, giving
\[
  \Prob(W^1 > 0) \;=\; \Prob(n \ge 2) \;=\; \sum_{n\ge2}(1-\rho)\rho^n
  \;=\; \rho^2,
  \qquad
  g^{\text{lateness}} \;=\; c_1\rho^2 .
\]

\paragraph{Tardiness (what the comment describes).}
Each job is charged once, at the moment it is found to be late.  With $D = 0$ a
job is late iff it does not enter service immediately, i.e.\ iff the server is
busy on arrival.  By PASTA the arriving customer sees the time-stationary
distribution, so
\[
  \Prob(\text{delay}) \;=\; \Prob(n \ge 1) \;=\; \sum_{n\ge1}(1-\rho)\rho^n \;=\;\rho ,
  \qquad
  g^{\text{tardiness}} \;=\; c\,\lambda\,\rho .
\]
The factor $\lambda$ is a unit conversion: $\Prob(\text{delay})$ is a fraction
\emph{of jobs}, and $g$ must be a cost \emph{per unit time}; multiplying by the
arrival rate converts one into the other.  This is also the bridge to the
service-level metric: with $A(T)$ arrivals and $L(T)$ of them late,
\[
  J \;=\; \lim_{T\to\infty}\frac{\E[L(T)]}{\E[A(T)]}
  \;=\; \frac{\lambda\rho\,T}{\lambda T} \;=\; \rho ,
\]
so the \emph{fraction} late is $\rho$ while the \emph{rate} of late jobs is
$\lambda\rho$; the two differ by the constant $\lambda$, and since $\lambda$ is
exogenous, minimising the rate is equivalent to minimising the fraction.

\paragraph{The degeneracy.}
The ratio of the two expressions is
\[
  \frac{g^{\text{tardiness}}}{g^{\text{lateness}}}
  \;=\; \frac{\lambda\rho}{\rho^2} \;=\; \frac{\lambda}{\rho} \;=\; \mu .
\]
\textbf{At $\mu = 1$ the two objectives predict exactly the same number.}
Table~2 uses $\mu = 1$ throughout, so agreement with $\rho^2$ confirms that the
cost normalisation is internally consistent but does \emph{not} confirm which
of the two cost functions is implemented --- precisely the distinction raised in
the introduction (see B.3).  Had tardiness been implemented by mistake,
Table~2 would look identical.

\paragraph{Correction.}
Repeat Experiment~1 with $\mu \neq 1$; e.g.\ $\mu = 2$, $\rho = 0.4$ gives
$\lambda = 0.8$ and the two predictions separate by a factor two:
$g^{\text{lateness}} = 0.16$ versus $g^{\text{tardiness}} = 0.32$.  This makes
the validation discriminating.  For $D > 0$ the tardiness benchmark also has a
closed form in M/M/1--FIFO, $\Prob(\text{wait} > D) = \rho e^{-\mu(1-\rho)D}$,
which reduces to $\rho$ at $D=0$ and provides a second, sharper test.

% ============================================================
\section*{B. Modelling questions answered by subsequent work}
% ============================================================

\subsection*{B.1 One event, several service starts --- the action space}

\comment{kan een event tot meer dan 1 service start leiden? indien max 1 dan
kies je de beste uit de toegelaten toewijzingen} \hfill (\S3.2)

\comment{als idleness mogelijk is zou een tick een trigger voor een service
start kunnen zijn} \hfill (\S3.2)

Yes, one event can trigger several service starts (a completion may free a
server while several types wait; a tick may make several assignments feasible
at once).  The manuscript resolves this by presenting candidates
\emph{sequentially} with binary accept/skip decisions.

\textbf{The alternative proposed in the comment --- ``choose the best among the
admissible assignments'' --- was subsequently implemented, and it resolves the
central failure of the manuscript's Experiment~2.}  In this
\emph{per-event} formulation each idle capacity unit makes a single
$(N{+}1)$-way decision: idle, or serve type $a \in \{1,\dots,N\}$, with
infeasible types masked.  Formally the two formulations admit the same physical
policies --- verified by three equivalence gates: the FIFO policy and the
never-serve policy reproduce bit-exactly, and the RVI optimum agrees on both
benchmark instances --- so what changes is only the \emph{credit structure}.

Why this matters is a property of the sequential decomposition that we have
called \emph{re-presentation forgiveness}.  A skipped job returns at the next
tick, so if a state is visited with slack for $m$ further presentations and the
policy assigns probability $p$ to serving, the job is served in time with
probability $1 - (1-p)^{m+1}$, which is close to $1$ even for $p$ well below
$\tfrac12$.  The stochastic policy is therefore behaviourally near-optimal even
where the sign of the logit is wrong; the advantage vanishes and training never
corrects it --- but the deployed policy is the \emph{argmax}, which turns such a
state into a deterministic error at every re-presentation.  In the per-event
space the same trade-off is a single logit comparison carrying the full
event-level advantage, and the pathology disappears.

Empirically, on Experiment~2 with the binary cost: the sequential space needs a
substantial stabilisation apparatus and reaches 6/8 seeds within 5\% of RVI; the
per-event space reaches \textbf{8/8 at the exact optimum with \emph{less}
machinery} (no temperature annealing).  The per-decision action set is also
$O(N)$ irrespective of the number of servers, which is the relevant property for
scaling.

Regarding the second comment: ticks do re-present the action queue in the
implementation, so a tick can indeed trigger a service start.  This is
deliberate, but it has a consequence worth stating in the manuscript: a higher
tick rate gives \emph{more} opportunities to reverse an idling decision, so
$\gamma$ is not entirely neutral with respect to the policy space.

\subsection*{B.2 Would queue lateness be better for reward shaping?}

\comment{is het vanuit het oogpunt van reward shaping niet beter queue lateness
te nemen?} \hfill (\S3)

This was tested directly.  The answer is nuanced: queue lateness does not solve
the shaping problem, and it introduces a harder one.

\paragraph{It does not fix sparsity.}
Queue lateness is denser than FIL lateness \emph{after} the deadline, but it is
equally flat \emph{before} it: both charge zero while $W_n^1 \le D$.  Since the
learning difficulty is the absence of gradient in the pre-deadline region, the
extension does not address it.

\paragraph{It introduces heavy tails.}
Let $B = \max(0, W - D)$.  As derived in C.4, the per-tick queue-lateness cost
is $\Theta(B^2)$.  Rare deep-backlog excursions therefore dominate the return
distribution: measured per-rollout reward rates vary from $-380$ to $-7500$
within otherwise healthy runs.  Under this variance the tuned recipe converges
to a \emph{bit-exact FIFO clone on 8 out of 8 seeds}.  Recovering optimal
behaviour required a different algorithmic configuration --- an average-reward
(differential) objective with a $64\times$ larger batch --- after which 8/8 seeds
reach the RVI optimum.

\paragraph{What does fix sparsity.}
A potential-based shaping term.  With $\Phi$ the accumulated urgency of the
waiting FIL, accruing $c_n/D$ per pre-deadline tick and refunded in full on
service, the shaped cost is $r' = r + \Phi(s) - \Phi(s')$.  This is
cost-neutral --- the potential telescopes, so every stationary policy retains its
long-run average, verified numerically for FIFO and RVI to four decimals --- while
supplying dense pre-deadline gradient.  It raises the success rate on the binary
cost from 2/8 to 6/8 seeds.

\paragraph{A dimensional caveat.}
The manuscript's normalisation by $\gamma$ (\S3.1) makes the FIL lateness cost
tick-rate invariant.  This does \emph{not} extend to queue lateness in its
excess-weighted form: measured FIFO cost is $1781 / 3035 / 5266$ at
$\gamma = 1.5 / 3 / 6$, i.e.\ approximately linear in $\gamma$, against
$74 / 71 / 70$ for the binary cost.  The reason is derived in C.4.

\subsection*{B.3 Lateness or tardiness?}

\comment{is dit niet tardiness?} \hfill (\S1)

\comment{dus je wijkt af van wat je in de intro zegt} \hfill (\S3.1)

The comments are justified: the introduction motivates the work by the fraction
of customers served late (tardiness), while the model optimises FIL lateness.
\S3 acknowledges the reason --- under tardiness, ``given fully flexible policy,
there is an incentive to never serve certain jobs once their FIL exceeds the due
date'' --- and defers the objective.

\textbf{That obstacle has since been removed.}  Two components were needed.

\paragraph{(i) An exact flux formulation.}
Charging at departure gives arbitrarily delayed credit.  Instead charge $c_n$ at
the instant a job crosses its deadline: the same total on every sample path, but
at the decision-relevant moment.  In the FIL projection this gives the
per-tick cost derived in C.5.

\paragraph{(ii) A rule that removes the degeneracy.}
The incentive identified in \S3 is real, and in the FIL projection it is worse
than stated: letting a FIL rot suppresses that type's arrival stream (arrivals
are only tracked when the queue is empty), so rotting is not merely free but
actively profitable.  This is an artefact of the projection, and RVI discovers
it: the computed optimum drifts below the achievable cost as the truncation
level grows, and learned policies ``beat'' the benchmark.

The remedy adopted is a modelling rule rather than a reward adjustment:
\emph{SLA escalation}.  When a capacity unit can serve a job already past its
deadline, the assignment is forced (oldest-late first, non-preemptive).  This is
interpretable as a service rule --- overdue work escalates to priority --- and is
implemented purely by restricting the admissible action set to a singleton, so
that PPO, DCL and RVI all solve the restricted MDP without any algorithmic
change.  Its effects are: the degeneracy disappears by construction; late states
become transient, so RVI truncation converges immediately (the computed optimum
is bit-identical at truncation levels $M = 30, 60, 90$, against a slow drift
without it); and the learned policy decides only \emph{pre-deadline}.

Under this objective the two-server benchmarks turn out to be trivially solvable
(FIFO and c-$\mu$ are both within $0.01\%$ of the optimum), and headroom appears
only on multi-hop routing structures.  This is a structural finding rather than
an algorithmic one, and it is the reason the current work has moved to designing
larger instances.

\subsection*{B.4 The FIL state is not sufficient --- what else can be done?}

\comment{eerst zeggen: \# bezette servers is niet voldoende info --- en ook:
huidige keuze niet optimaal. verder nog iets mee doen?} \hfill (\S2.2)

Agreed on both counts, and the manuscript should say so explicitly.  The
implementation already supports a generalisation: the number of tracked queue
positions per type is a parameter $d$ (\texttt{max\_queue\_depth}), with the
untracked remainder handled by the conditional-expectation tail derived in C.4;
at $d = 1$ it reduces exactly to the FIL model.  The observation supplied to the
network is a separate parameter, so simulator fidelity and policy input can be
varied independently.  This gives a natural sensitivity analysis --- how much does
the FIL projection cost? --- which has not yet been run and would strengthen the
paper.

One point should be stated precisely, since it recurs below: the projected
\emph{cost} is exact (see C.4), whereas the projected \emph{dynamics} are an
approximation, because the untracked queue influences future transitions.  RVI,
the simulator and the critic all use the same projection, so comparisons are
internally consistent; but the optimum being computed is that of the projected
model.

% ============================================================
\section*{C. Clarifications}
% ============================================================

\subsection*{C.1 Why truncation impedes the span criterion}

\comment{dit snap ik niet: waarom zou deze afknotting de convergentie
verhinderen?} \hfill (\S4)

The manuscript's phrasing is too strong and should be weakened.  Two distinct
convergence questions are involved.

\paragraph{(i) Iteration at fixed $M$.}
For a unichain aperiodic model the span seminorm
$\mathrm{sp}(h_{t+1} - h_t)$ does converge to zero; truncation does not
\emph{prevent} this.  What it does is slow it down: clamping at $W = M$ creates
boundary states with a self-loop (a job at level $M$ remains at $M$ with high
probability), and this slowly mixing region keeps a small residual spread alive
for many iterations after $g_t$ has stabilised.  Since $g_t$ is an average it
converges much faster.  The accurate statement is therefore that the span is
\emph{too slow to be a practical stopping criterion within the iteration
budget}, not that it never converges --- and the two-stage criterion is a
pragmatic response to that.

\paragraph{(ii) Convergence in $M$.}
This is the more consequential issue and is not discussed in the manuscript.
The truncation level is increased until the relative change in $g^*$ between
successive levels falls below $1\%$.  For the binary cost this is adequate, since
the cost saturates in $W$.  For tail-heavy objectives it is not: the cost keeps
growing with queue depth, so $g^*$ still drifts at the stopping point.
Tightening the tolerance to $10^{-3}$ moved the queue-lateness reference on
Experiment~2 from $2492$ to $2350$, a $5.7\%$ correction.  This was discovered
because \emph{learned policies evaluated below the computed optimum} --- a useful
diagnostic worth reporting: any policy beating the benchmark indicates a loose
benchmark, not a good policy.

\subsection*{C.2 DCL versus standard deep RL}

\comment{wat is het verschil tussen DCL en ``standaard'' deep RL? en waarom dan
DCL?} \hfill (\S5.1)

\comment{als ik het goed begrijp heb je alleen een NN om de policy te
beschrijven, niet voor de waardefunctie} \hfill (Alg.~1)

\comment{gebruik je verdiscontering?} \hfill (Alg.~1)

\comment{zijn er convergentiegaranties met DCL?} \hfill (\S5.2)

The second comment is correct: DCL trains a network for the policy only.  The
distinctions are as follows.

\begin{center}
\begin{tabular}{lll}
\toprule
 & \textbf{DCL} & \textbf{PPO (policy gradient)} \\
\midrule
Networks & policy only & policy (actor) + value (critic) \\
Learning signal & best action per state, by rollouts & advantage estimates \\
Loss & cross-entropy (classification) & clipped surrogate + value loss \\
Discounting & none; undiscounted horizon $H$ & $\gamma^{\Delta}$, or average-reward \\
Variance control & CRN + sequential halving & GAE, batching, normalisation \\
\bottomrule
\end{tabular}
\end{center}

DCL replaces value regression by \emph{best-action identification}: it does not
estimate how good actions are, only which is best, which is a statistically
easier question and is why CRN and sequential halving are effective.  The price
is that policy improvement is only as good as the labels, and that no value
function is available for bootstrapping.

\paragraph{Convergence guarantees.}
None in this setting.  DCL is approximate policy iteration with classification
error, and the standard guarantees require the classification error to be
controlled uniformly, which is not verifiable here.  Empirically the failure
mode is a \emph{fixed point}: started from a deterministic FIFO base, DCL
frequently returns an exact FIFO clone, because on FIFO's own state
distribution the deviations do not appear advantageous, and generation~2 starts
from the same policy.  Widening the base distribution (a stochastic base
policy) is a partial remedy.

The same caveat applies to PPO: with neural function approximation there is no
convergence proof.  What replaces guarantees in this work is (i) an exact
benchmark on instances small enough for RVI, and (ii) reporting every claim as
a seed distribution --- ``$k$ of 8 seeds within $x\%$ of RVI'' --- rather than a
single run.

\subsection*{C.3 What the value function is, in the PPO setting}

Since the manuscript version predates the PPO work, this is background rather
than a response, but it is the natural follow-up to the previous comment.

PPO maintains a critic $V_\phi$ alongside the policy.  Under the discounted
objective it approximates $V^\pi_\gamma$.  Under the average-reward
(differential) objective --- which is what the queue-lateness configuration
uses --- it approximates the \emph{differential value function} $h^\pi$ solving
the Poisson equation
\[
  h^\pi(x) \;=\; r(x) - g^\pi \;+\; \sum_{x'} P^\pi(x'\mid x)\, h^\pi(x') ,
\]
with $g^\pi$ estimated online.  \textbf{This is the same object that RVI
computes}: RVI solves it exactly by sweeps on a truncated state space, the
critic approximates it by regression on simulated returns over the full space.
That correspondence is what makes RVI a meaningful benchmark rather than merely
a convenient one.

Three implementation points are load-bearing and each was validated by ablation
(8 seeds; the full recipe attains 6/8):
\begin{itemize}
  \item \textbf{Semi-MDP discounting.}  Consecutive decisions are separated by
  a variable number $\Delta_t$ of uniformised periods, including
  $\Delta_t = 0$ for two decisions within one tick.  The discount must be
  $\gamma^{\Delta_t}$ with $\gamma^0 = 1$.  Using $\gamma^{\max(1,\Delta_t)}$
  instead subsidises idling --- most strongly in busy states, where the
  candidate list is longest --- and yields \textbf{0/8} successful seeds, with
  failures at the never-serve policy.
  \item \textbf{Value-target normalisation.}  Returns are $O(10^2$--$10^3)$
  while the policy gradient operates on standardised advantages of order one.
  Without normalising the value targets the value loss dominates the shared
  trunk: \textbf{0/8}, and perfectly bimodal --- every seed converges either to
  the exact never-serve policy or to an exact FIFO clone.  Normalisation does
  not change the fixed point, only the conditioning of the optimisation.
  \item \textbf{Staggered environment resets.}  With persistent environments a
  deep-backlog excursion moves all data into a region where the shaping term
  has saturated; the policy's behaviour on the states it will actually be
  deployed in then degrades unobserved.  Resetting one environment per update
  gives \textbf{1/8} $\to$ 6/8.
\end{itemize}

\subsection*{C.4 The queue-lateness cost, derived}

This addresses \comment{hoe bereken je de kosten exact?} for the full-queue
variant, and supplies the derivation underlying B.2.

Condition on the FIL of type $n$ having level $W = w$.  Under FIFO within a
type, no job that arrived after the FIL can have been served, so the jobs behind
the FIL are exactly the type-$n$ arrivals of the last $w$ ticks; index the
tick-slots $i = 0,\dots,w-1$ by the age of a job arriving in them.

In the uniformised chain, arrivals and ticks compete, so the number $K$ of
arrivals in one slot is geometric,
\[
  \Prob(K = k) \;=\; \alpha^k\beta,
  \qquad \alpha = \frac{\lambda}{\lambda+\gamma},
  \quad \beta = \frac{\gamma}{\lambda+\gamma},
  \qquad \E[K] \;=\; \frac{\alpha}{\beta} \;=\; \frac{\lambda}{\gamma},
\]
consistent with continuous time ($\lambda/\gamma$ expected arrivals per slot of
length $1/\gamma$).  Note this is exactly the model encoded by the FIL-refresh
law~(3): $\Prob(K = 0) = \beta$ is the probability that a slot is empty.

\paragraph{Two distinct quantities.}
$\alpha$ is the probability that a slot is \emph{non-empty}; $\lambda/\gamma$ is
the \emph{expected number} of jobs in it.  Which one is required depends on
whether the cost counts jobs or detects presence --- a distinction that has
practical consequences (C.5).

With $B = \max(0, W-D)$, summing the excess over the untracked jobs:
\[
  \E\Big[\sum_{j\text{ behind}} \max(0, A_j - D)\Big]
  \;=\; \sum_{i=0}^{w-1}\E[K]\max(0,i-D)
  \;=\; \frac{\lambda}{\gamma}\sum_{j=1}^{B-1} j
  \;=\; \frac{\lambda}{\gamma}\cdot\frac{(B-1)B}{2}.
\]
The same result follows from the conditional-uniformity property invoked in the
manuscript: given $q$ arrivals in the window their ages are distributed as the
order statistics of $q$ uniforms, so the expectation factorises as
$\E[q]\cdot\E[\max(0,U-D)] = \frac{\lambda}{\gamma}w \cdot \frac{1}{w}\frac{(B-1)B}{2}$.
Substituting the expectation is exact because the cost is \emph{linear} in the
individual excesses; there is no Jensen gap.  Hence
\[
  \boxed{\;r_{\text{QL}} \;=\; \frac{c}{\gamma}\left[B \;+\;
  \frac{\lambda}{\gamma}\cdot\frac{(B-1)B}{2}\right]\;}
\]
which generalises to tracking depth $d>1$ by charging the tracked positions
exactly and anchoring the tail at the deepest tracked position.

\paragraph{A discrepancy with Table~1.}
The manuscript defines FIL lateness with an \emph{indicator},
$\frac{c}{\gamma}\mathbf{1}\{W > D\}$, so its natural full-queue extension is a
\emph{count} of overdue jobs, with tail $\frac{\lambda}{\gamma}(B-1)^+$.  The
implementation instead weights by the \emph{excess}, as above.  These are
different objectives: a job overdue for $T$ ticks accumulates $\Theta(T)$ under
the count version but $\sum_{k\le T}k = \Theta(T^2)$ under the excess version.
Both are defensible, but they must be named correctly --- and the choice explains
two observations at once:

\begin{center}
\begin{tabular}{lcc}
\toprule
variant & cost per unit time & tick-rate invariant? \\
\midrule
count  & $c \cdot \#\{\text{overdue}\}$ & yes \\
excess & $c\,\gamma \sum_j \tau_j$ & no, scales with $\gamma$ \\
\bottomrule
\end{tabular}
\end{center}
Here $\tau_j$ is overdue time in physical units, and excess in ticks is
$\tau\gamma$.  The excess variant is therefore both the source of the heavy
tails (B.2) and of the tick-rate dependence, and an additional factor $1/\gamma$
restores invariance.

\subsection*{C.5 The tardiness flux cost, derived}

For completeness, the objective discussed in B.3.  Crossings occur when a job's
age passes from $D$ to $D+1$:
\begin{itemize}
  \item the FIL crosses exactly when $W = D+1$ after the tick;
  \item jobs behind it occupy ages $0,\dots,W-1$, so the slot reaching $D+1$
  lies strictly behind the FIL only when $W \ge D+2$; the expected number of
  jobs in that slot is $\E[K] = \lambda/\gamma$.
\end{itemize}
\[
  \boxed{\;r_{\text{QT}} \;=\; c\left[\mathbf{1}\{W = D+1\}
  \;+\; \frac{\lambda}{\gamma}\,\mathbf{1}\{W \ge D+2\}\right]\;}
\]
Here the coefficient must be the expected \emph{number} $\lambda/\gamma$ and not
the occupancy probability $\alpha$, since a slot containing $k$ jobs produces
$k$ crossings.  The distinction is verifiable: if the queue stays backed up,
every arriving job eventually crosses, so crossings per unit time must equal
$\lambda$; and indeed $\gamma \cdot (\lambda/\gamma) = \lambda$ for every
$\gamma$, whereas $\gamma\alpha = \lambda\beta$ depends on $\gamma$.  This makes
the objective tick-rate invariant up to $O(1/\gamma)$ discretisation terms --- the
crossing is registered at age $D + 1/\gamma$ rather than $D$, and the guard
$W\ge D+2$ is $2/\gamma$ beyond the deadline; both vanish as $\gamma$ grows.

% ============================================================
\section*{D. Terminology and notation}
% ============================================================

\begin{center}
\begin{tabular}{p{7.4cm}p{7.4cm}}
\toprule
\textbf{Comment} & \textbf{Response} \\
\midrule
\comment{deadline $d_n$? een vaste ``acceptable waiting time'' afh van het
jobtype?} / \comment{het is niet echt een deadline, dat is de tijd waarop de
job in bediening moet/klaar moet zijn}
& Agreed.  $D_n$ is a type-dependent \emph{acceptable waiting time until
service start}.  The manuscript should state ``service start'' explicitly, as
the distinction from service completion is currently only implicit. \\[4pt]
\comment{$N$ heb je al eerder gebruikt (\# job types)}
& Trajectories per generation renamed to $T$. \\[4pt]
\comment{$M$ heb je al 2 keer eerder gebruikt (voor \# server en voor
afknottingsniveau)}
& Truncation level renamed $\bar{W}$, simulation budget $B$; $M_k$ retained for
the number of servers in group $k$. \\[4pt]
\comment{ik geef de voorkeur aan wiskundige notatie: $=$ ipv $\leftarrow$,
$\mathbb{P}$ (of blackboard P) ipv Pr}
& Adopted throughout. \\[4pt]
\comment{shaded $\to$ bold} (Table 1)
& Adopted. \\[4pt]
\comment{referentie} (Adam)
& Kingma \& Ba (2015) added. \\[4pt]
\comment{offers little room for suboptimal data}
& Rephrased: in the single-server M/M/1 case the FIFO base policy already
visits essentially the whole relevant state space, so the sampled data contain
few informative deviations. \\[4pt]
Repeated ``Hence'' (\S6.1)
& Rewritten. \\
\bottomrule
\end{tabular}
\end{center}

\end{document}
